\section{Scoring formulas}
\label{sec:formulas}

\subsection{Notation}

All scores are on a \numrange{0}{100} scale unless stated otherwise. Two
bounding operators are used throughout:

\begin{align}
\clamp{x}          &= \min\bigl(100,\; \max(0,\; x)\bigr), \\
\clip{x}{l}{u}     &= \min\bigl(u,\; \max(l,\; x)\bigr).
\end{align}

$\clampop$ bounds a score to the reporting scale; $\clipop$ bounds a physical
estimate to a stated envelope. The distinction matters in
\Cref{sec:cooling}, where a score is clamped to $[0,100]$ while a temperature
is clipped to a literature-derived interval.

Every weight appearing below is reproduced with its rationale in
\Cref{app:weights}. Within each weighted sum, the weights sum to unity.

\subsection{Heat Exposure Score}
\label{sec:heat-exposure}

There are two paths, selected by data availability.

\paragraph{Data-rich path.} Where a land-surface-temperature anomaly is
supplied and converted to $S_{\mathrm{LST}}$ by \Cref{eq:lst}:

\begin{equation}
\label{eq:heat-exposure-rich}
S_{\mathrm{heat}} \;=\; \clampop\bigl(
  0.40\,S_{\mathrm{LST}}
+ 0.25\,S_{\mathrm{imp}}
+ 0.20\,S_{\mathrm{solar}}
+ 0.15\,S_{\mathrm{vegdef}}\bigr),
\end{equation}

where $S_{\mathrm{imp}}$ is the impervious surface percentage of the site used
directly as a \numrange{0}{100} score, $S_{\mathrm{solar}}$ is the qualitative
solar exposure level normalised by \Cref{tab:levels}, and the vegetation
deficit is

\begin{equation}
\label{eq:vegdef}
S_{\mathrm{vegdef}} \;=\; \clamp{100 - g},
\end{equation}

with $g$ the existing green cover percentage.

\paragraph{Data-poor path.} Where no LST anomaly is supplied, the Heat Exposure
Score is the normalised qualitative heat exposure level directly:
$S_{\mathrm{heat}} = S_{\mathrm{heatlevel}}$.

\paragraph{Basis.} The component selection reflects established
surface-energy-balance drivers of urban heat: surface temperature, sealed
surfaces, solar loading, and vegetation deficit. \textcite{ziter2019} provides
direct empirical support for two of the four --- impervious cover raises air
temperature approximately linearly, and canopy cover reduces it nonlinearly ---
and additionally establishes that the warming effect of impervious surfaces
persists at night while canopy cooling largely does not.

\paragraph{What is expert judgment here.} The \emph{relative} weights are not
derived from literature. No retrieved source ranks the marginal contribution of
LST anomaly against imperviousness against solar loading against vegetation
deficit for the purpose of a screening index. The weights express a judgment
that a directly measured thermal signal should dominate where it is available,
that sealed surface is the strongest single proxy where it is not, and that
vegetation deficit --- already partly captured by the other three --- should
carry the least. This judgment is declared, versioned, and subject to the
sensitivity analysis of \Cref{sec:sensitivity}.

\subsection{Vulnerability Score}
\label{sec:vulnerability}

\begin{equation}
\label{eq:vulnerability}
S_{\mathrm{vuln}} \;=\; \clampop\bigl(
  0.40\,S_{\mathrm{dens}}
+ 0.40\,S_{\mathrm{groups}}
+ 0.20\,S_{\mathrm{coolaccess}}\bigr),
\end{equation}

where $S_{\mathrm{dens}}$ is normalised population density,
$S_{\mathrm{groups}}$ is normalised vulnerable-group presence, and
$S_{\mathrm{coolaccess}}$ is the cooling access \emph{deficit} obtained through
the inverted mapping of \Cref{sec:inverted}.

\paragraph{Basis.} Indicator selection follows \textcite{reid2009}, whose
factor analysis of ten heat-vulnerability variables identified four factors
explaining more than \SI{75}{\percent} of total variance: social and
environmental vulnerability, social isolation, air-conditioning prevalence, and
the proportion of elderly and diabetic residents.

Population density and vulnerable-group presence are weighted \emph{equally}
because that analysis provides no basis for ranking demographic vulnerability
above exposure density; asserting a difference would be inventing a finding.
\textcite{sera2019} independently supports density as a modifier of the
heat--mortality relationship across \num{340} cities in \num{22} countries.
Cooling access carries the smaller weight because it is the field most
frequently unknown at screening stage, so a larger weight would amplify the
neutral default of \Cref{sec:unknown}.

\paragraph{One dimension, two indicators.} From version \texttt{2026.08.03}
the cooling access deficit is measured by two indicators rather than one, at
the same \num{0.20} weight: access to actively cooled \emph{indoor} space, and
access to a cool \emph{outdoor} refuge --- shaded green or blue space the
site's users can reach. The split is evidenced, not assumed.
\textcite{reid2009} grounds the indoor indicator. \textcite{burkart2016}
grounds the outdoor one directly: in Lisbon, above the 99th temperature
percentile, mortality among residents over \num{65} rose
\SI{14.7}{\percent} per \si{\degreeCelsius} in the least-vegetated areas
against \SI{3.0}{\percent} in the most vegetated, and \SI{7.1}{\percent} per
\si{\degreeCelsius} beyond \SI{4}{\kilo\metre} from water against
\SI{2.1}{\percent} within it.

The dimension is the \emph{mean of the indicators actually supplied}:

\begin{equation}
\label{eq:coolaccess}
S_{\mathrm{coolaccess}} \;=\;
\begin{cases}
  \dfrac{1}{\lvert \mathcal{A}\rvert} \displaystyle\sum_{i \in \mathcal{A}} S_i
    & \text{if } \mathcal{A} \neq \emptyset,\\[2ex]
  50 & \text{otherwise,}
\end{cases}
\end{equation}

where $\mathcal{A}$ is the set of supplied indicators (an explicit
\emph{unknown} is not supplied). Averaging an \emph{unsupplied} indicator in at
the neutral \num{50} was considered and rejected: it would pull a signal the
user did give toward the middle purely because a second question was skipped,
letting an unknown alter the meaning of an answer --- the defect removed from
the soil--water pair in \Cref{sec:adjustment}. Under the adopted rule a single
answer speaks for the dimension, answering the second moves it in whichever
direction that answer warrants, and skipping costs nothing. The two indicators
correspondingly share one confidence slot (\Cref{sec:confidence}).

\paragraph{A disclosed overlap.} The outdoor indicator describes cool refuges
in the \emph{surroundings} that the site's users can reach; the vegetation
deficit of \Cref{sec:heat-exposure} describes greenness \emph{of the site
itself}. The two are distinct quantities but conceptually adjacent, and a bare
site in a green district registers a low outdoor-refuge deficit alongside a
high vegetation deficit. We state this rather than leave it for a reader to
discover, as with the vulnerability double contribution of
\Cref{sec:aggregation}.

\subsection{Heat Priority Index}
\label{sec:hpi}

\begin{equation}
\label{eq:hpi}
\mathrm{HPI} \;=\; \clampop\bigl(0.60\,S_{\mathrm{heat}} + 0.40\,S_{\mathrm{vuln}}\bigr).
\end{equation}

The index is interpreted by the bands in \Cref{tab:hpi-bands}. Band boundaries
drive recommendation text, which is why their stability under weight
perturbation is one of the four quantities the sensitivity analysis reports
(\Cref{sec:sensitivity-results}).

\begin{table}[H]
\centering
\caption{Interpretation bands for the Heat Priority Index and the NbS Cooling
Opportunity Score.}
\label{tab:hpi-bands}
\small
\begin{tabular}{@{}lll@{}}
\toprule
\textbf{Range} & \textbf{Heat Priority Index} & \textbf{Opportunity Score} \\
\midrule
0\,--\,30   & Low      & Low priority \\
31\,--\,60  & Medium   & Moderate opportunity \\
61\,--\,80  & High     & Strong opportunity \\
81\,--\,100 & Critical & High-priority project \\
\bottomrule
\end{tabular}
\end{table}

\subsection{Derived site adjustment factor}
\label{sec:adjustment}

Four site conditions modulate typology performance. Each condition resolves to
a level, and each level maps to a multiplicative factor:

\begin{center}
\small
\begin{tabular}{@{}lccccc@{}}
\toprule
\textbf{Condition level} & Poor & Moderate & Good & Excellent & Unknown \\
\midrule
\textbf{Factor} & 0.5 & 0.8 & 1.0 & 1.2 & 1.0 \\
\bottomrule
\end{tabular}
\end{center}

The composite adjustment factor is

\begin{equation}
\label{eq:adjustment}
A \;=\; 0.40\,f_{\mathrm{canopy}}
      + 0.25\,f_{\mathrm{soilwater}}
      + 0.20\,f_{\mathrm{scale}}
      + 0.15\,f_{\mathrm{climate}},
\end{equation}

which ranges over $[0.5, 1.2]$, attaining its bounds only when all four
conditions are poor or all four are excellent.

\paragraph{All four conditions are derived, not asked.} No additional questions
are put to the user; each condition is computed from inputs already supplied.
\Cref{tab:derivation} gives the rules.

\begin{table}[htbp]
\centering
\caption[Derivation of site adjustment conditions]{Derivation of the four site
adjustment conditions from inputs the user has already supplied.}
\label{tab:derivation}
\small
\begin{tabularx}{\textwidth}{@{}l>{\raggedright\arraybackslash}X@{}}
\toprule
\textbf{Condition} & \textbf{Derivation rule} \\
\midrule
Canopy &
$c = \dfrac{\text{existing tree canopy} + \text{new canopy at maturity}}
           {\text{site area}}\times 100$, then:
poor if $c < 10$; moderate if $10 \le c < 25$; good if $25 \le c < 40$;
excellent if $c \ge 40$. \\
\addlinespace
Soil--water &
Soil availability and irrigation availability are each mapped to a condition
level, and the \emph{more limiting} of the two is taken. The mapping is
\emph{none}\,$\rightarrow$\,poor;
\emph{limited}, \emph{occasional}\,$\rightarrow$\,moderate;
\emph{moderate}\,$\rightarrow$\,good;
\emph{reliable}, \emph{high}\,$\rightarrow$\,excellent.
When exactly one of the two is known, the condition is capped at \emph{good}:
a pair containing an unknown may never exceed the neutral factor, so declaring
reliable irrigation is never worse --- and with high soil strictly better ---
than skipping the question (D-026, version \methver{}). \\
\addlinespace
Scale &
From the declared assessment scale:
\emph{city}, \emph{district}\,$\rightarrow$\,excellent;
\emph{neighbourhood}\,$\rightarrow$\,good;
\emph{site}, \emph{building}\,$\rightarrow$\,moderate. \\
\addlinespace
Climate &
A lookup of typology family against climate zone
(\Cref{tab:climate-matrix}). \\
\bottomrule
\end{tabularx}
\end{table}

\paragraph{Why canopy carries the largest weight.} It is the condition with the
strongest direct empirical support. \textcite{ziter2019} found that cooling
increases nonlinearly with canopy cover and is greatest above approximately
\SI{40}{\percent} cover; the \emph{excellent} boundary is placed at that
empirically identified inflection rather than chosen for convenience. The same
study locates maximum daytime cooling at \SIrange{60}{90}{\metre} radii, which
is the basis for the scale condition: a district-scale or city-scale
intervention is more likely to reach the spatial scale at which the effect is
strongest than a single-building one.

\paragraph{Why climate is present at all.} \textcite{keravec2026} finds that
adding K\"oppen--Geiger subzone raises explained variance in cooling
effectiveness from \SI{12}{\percent} to \SI{78}{\percent}. Omitting climate
would discard the largest single explanatory term in the available evidence.

\begin{table}[htbp]
\centering
\caption[Climate suitability matrix]{Climate suitability matrix: the condition
level assigned to each typology family in each supported climate zone, resolved
to a factor through the condition-level table above. Vegetation-based
typologies are downgraded in arid and semi-arid zones where water limitation
constrains evapotranspiration; blue-green typologies are downgraded in
tropical-wet zones following the humidity caveat of
\textcite{gunawardena2017}, and upgraded in tropical-dry zones. An
unrecognised zone yields \emph{unknown}, whose factor is the neutral 1.0.}
\label{tab:climate-matrix}
\small
\begin{tabular}{@{}lllll@{}}
\toprule
\textbf{Climate zone} & \textbf{Green} & \textbf{Building} &
\textbf{Blue-green} & \textbf{Hybrid} \\
\midrule
Tropical wet & Good     & Good     & Moderate  & Good \\
Tropical dry & Good     & Good     & Excellent & Good \\
Arid         & Moderate & Moderate & Moderate  & Good \\
Semi-arid    & Moderate & Good     & Good      & Good \\
Temperate    & Good     & Good     & Good      & Good \\
Other        & Unknown  & Unknown  & Unknown   & Unknown \\
\bottomrule
\end{tabular}
\end{table}

\subsection{Cooling Potential Score and temperature reduction}
\label{sec:cooling}

The Cooling Potential Score scales the typology's base cooling score by the
site adjustment factor:

\begin{equation}
\label{eq:cps}
S_{\mathrm{cool}} \;=\; \clamp{B \times A},
\end{equation}

where $B$ is the base cooling score from \Cref{tab:typologies}.

The reported temperature reduction range is computed differently, and the
difference is the point of this subsection:

\begin{align}
\label{eq:dtmin}
\dT_{\min} &= \clip{\tau_{\min} \times A}{\tau_{\min}}{\tau_{\max}}, \\
\label{eq:dtmax}
\dT_{\max} &= \clip{\tau_{\max} \times A}{\tau_{\min}}{\tau_{\max}},
\end{align}

where $[\tau_{\min}, \tau_{\max}]$ is the typology's literature envelope from
\Cref{tab:typologies}.

\begin{keynote}
\noindent\textbf{The clipping rule is a deliberate correction to the draft
methodology.} Without it, an adjustment factor above 1.0 applied to a
typology's upper bound produces temperature claims exceeding anything in the
retrieved literature. An urban forest on an excellent site
($A = 1.2$, $\tau_{\max} = \qty{3.0}{\degreeCelsius}$) would report up to
\qty{3.6}{\degreeCelsius} --- a figure no retrieved source supports for a
site-level intervention.

Because the published envelopes already represent well-executed interventions
under favourable conditions, favourable site conditions must not be permitted
to compound that favourability. The \numrange{0}{100} Cooling Potential Score
of \Cref{eq:cps} remains free to rise above the base value, so the tool can
still express that a site is exceptionally well suited --- but it expresses
this as a higher relative score, not as a larger physical claim.
\end{keynote}

\Cref{fig:clipping} shows the behaviour across the full range of the adjustment
factor. The rule is asymmetric in effect, and it is worth being precise about
what it does. For $A < 1$ the upper bound falls --- a poor site is reported as
delivering less --- while the lower bound is held at $\tau_{\min}$, so the
reported range narrows from the top. For $A > 1$ the upper bound is held at
$\tau_{\max}$ while the lower bound rises, so the range narrows from the
bottom. In both directions the reported interval remains a subinterval of the
literature envelope: site conditions shift an estimate \emph{within} the
published evidence, and never outside it.

\begin{figure}[htbp]
\centering
% The temperature-envelope clipping rule (decision D-008).
% Illustrated for street tree planting: envelope 0.5-3.0 C, with the site
% adjustment factor ranging over its full 0.5-1.2 span. The unclipped product
% is shown dashed; the reported values are the clipped ones. The vertical
% marker locates the worked example of Section 13.
\begin{tikzpicture}
\begin{axis}[
  font=\footnotesize,
  width=0.80\textwidth,
  height=74mm,
  xmin=0.5, xmax=1.22,
  ymin=0, ymax=3.9,
  xtick={0.5,0.6,...,1.2},
  xticklabel style={/pgf/number format/fixed, /pgf/number format/precision=1},
  ytick={0,0.5,...,3.5},
  yticklabel style={/pgf/number format/fixed, /pgf/number format/precision=1},
  xlabel={Site adjustment factor $A$},
  ylabel={Reported temperature reduction (\si{\degreeCelsius})},
  xlabel style={font=\footnotesize},
  ylabel style={font=\footnotesize},
  axis lines=left,
  axis line style={draw=rulegrey, line width=0.5pt},
  tick style={draw=rulegrey},
  grid=both,
  grid style={draw=rulegrey!40, line width=0.3pt, dash pattern=on 1pt off 1.5pt},
  legend style={at={(0.5,-0.20)}, anchor=north, font=\scriptsize,
    legend columns=1, draw=rulegrey, fill=white, legend cell align=left,
    row sep=1.5pt},
  clip=true,
]

% --- literature envelope band -------------------------------------------
\addplot[draw=none, fill=criterragreen!12, forget plot]
  coordinates {(0.5,0.5) (1.22,0.5) (1.22,3.0) (0.5,3.0)} \closedcycle;

% --- unclipped upper bound: 3.0 * A --------------------------------------
\addplot[dashed, draw=criterraink, line width=0.9pt, domain=0.5:1.22, samples=2]
  {3.0*x};
\addlegendentry{Unclipped $\Delta T_{\max}=3.0\,A$ --- rejected as unsupported}

% --- clipped upper bound: clip(3.0*A, 0.5, 3.0) --------------------------
\addplot[draw=criterragreen, line width=1.5pt, domain=0.5:1.0, samples=2] {3.0*x};
\addlegendentry{Reported $\Delta T_{\max}=\mathrm{clip}(3.0\,A,\;0.5,\;3.0)$}
\addplot[draw=criterragreen, line width=1.5pt, domain=1.0:1.22, samples=2,
  forget plot] {3.0};

% --- clipped lower bound: clip(0.5*A, 0.5, 3.0) --------------------------
\addplot[draw=criterragreen, line width=1.5pt, dash pattern=on 3pt off 2pt,
  domain=0.5:1.0, samples=2] {0.5};
\addlegendentry{Reported $\Delta T_{\min}=\mathrm{clip}(0.5\,A,\;0.5,\;3.0)$}
\addplot[draw=criterragreen, line width=1.5pt, dash pattern=on 3pt off 2pt,
  domain=1.0:1.22, samples=2, forget plot] {0.5*x};

% --- envelope boundary guides -------------------------------------------
\addplot[draw=rulegrey, line width=0.5pt, domain=0.5:1.22, samples=2,
  forget plot] {3.0};
\node[anchor=south west, font=\scriptsize\itshape, color=criterraink,
      inner sep=1.5pt] at (axis cs:0.51,3.01)
  {literature ceiling \SI{3.0}{\degreeCelsius}};
\node[anchor=south west, font=\scriptsize\itshape, color=criterraink,
      inner sep=1.5pt] at (axis cs:0.51,0.51)
  {literature floor \SI{0.5}{\degreeCelsius}};

% --- suppressed inflation -----------------------------------------------
\draw[{Latex[length=1.8mm]}-{Latex[length=1.8mm]}, draw=criterraink,
  line width=0.6pt] (axis cs:1.20,3.02) -- (axis cs:1.20,3.58);
\node[anchor=east, font=\scriptsize\itshape, align=right, color=criterraink,
      inner sep=2pt] at (axis cs:1.195,2.68)
  {clipping suppresses\\this inflation};
\draw[draw=criterraink, line width=0.5pt] (axis cs:1.20,2.86) -- (axis cs:1.20,3.00);

% --- the worked example of Section 13 -----------------------------------
\draw[dotted, draw=criterraink, line width=0.8pt]
  (axis cs:0.95,0) -- (axis cs:0.95,2.85);
\addplot[only marks, mark=*, mark size=2pt,
  mark options={draw=criterraink, fill=white, line width=0.8pt}, forget plot]
  coordinates {(0.95,2.85) (0.95,0.5)};
\node[anchor=south east, font=\scriptsize\itshape, align=right,
      color=criterraink, inner sep=2pt] at (axis cs:0.945,2.86)
  {worked example:\\$A=0.95$};

\end{axis}
\end{tikzpicture}

\caption[The temperature-envelope clipping rule]{The temperature-envelope
clipping rule of \Cref{eq:dtmin,eq:dtmax}, illustrated for street tree planting
(envelope \SIrange{0.5}{3.0}{\degreeCelsius}) across the full
\numrange{0.5}{1.2} range of the site adjustment factor $A$. The dashed line is
the unclipped product $3.0\,A$, which the methodology rejects because for
$A > 1$ it exceeds the literature ceiling. The reported bounds remain within
the shaded envelope throughout. The marked point locates the worked example of
\Cref{sec:worked-example}.}
\label{fig:clipping}
\end{figure}

\paragraph{Derived reporting quantities.} Two further quantities are reported
directly. Shade potential is the fraction of the site brought under new canopy
at maturity:

\begin{equation}
\label{eq:shade}
S_{\mathrm{shade}} \;=\; \clamp{\frac{a_{\mathrm{canopy}}}{a_{\mathrm{site}}}\times 100},
\end{equation}

with $a_{\mathrm{canopy}}$ the new canopy area at maturity and
$a_{\mathrm{site}}$ the site area, both in square metres. And the midpoint of
the adjusted temperature range,
$\overline{\dT} = (\dT_{\min} + \dT_{\max})/2$, is assigned a heat-index
improvement bucket: \emph{low} below \qty{0.5}{\degreeCelsius}, \emph{medium}
from \qty{0.5}{\degreeCelsius} to \qty{1.5}{\degreeCelsius}, and \emph{high}
from \qty{1.5}{\degreeCelsius} upward.

\paragraph{Time to benefit.} No maturity discount is applied to cooling values.
Instead the expected maturity period is reported as a time-to-benefit class, so
that a slow-maturing intervention is never presented as delivering its cooling
immediately: under one year \emph{immediate}, one to three years
\emph{short-term}, three to ten years \emph{medium-term}, ten years and beyond
\emph{long-term} (half-open brackets). Absent a maturity period the output is
\emph{not estimated}.

\subsection{NbS Suitability, co-benefit, and equity scores}
\label{sec:composite-subscores}

Three further composite sub-scores enter the result. Their aggregation weights
are declared in configuration and reproduced in \Cref{app:weights}:

\begin{align}
\label{eq:suitability}
S_{\mathrm{suit}} &= \clampop\bigl(
  0.30\,S_{\mathrm{space}} + 0.25\,S_{\mathrm{soil}} + 0.20\,S_{\mathrm{water}}
+ 0.15\,S_{\mathrm{maint}} + 0.10\,S_{\mathrm{context}}\bigr), \\[2pt]
\label{eq:cobenefit}
S_{\mathrm{cobenefit}} &= \clampop\bigl(
  0.25\,S_{\mathrm{biodiv}} + 0.25\,S_{\mathrm{storm}} + 0.20\,S_{\mathrm{health}}
+ 0.15\,S_{\mathrm{social}} + 0.15\,S_{\mathrm{urbanq}}\bigr), \\[2pt]
\label{eq:equity}
S_{\mathrm{equity}} &= \clampop\bigl(
  0.35\,S_{\mathrm{vulnben}} + 0.25\,S_{\mathrm{access}} + 0.20\,S_{\mathrm{safety}}
+ 0.20\,S_{\mathrm{particip}}\bigr).
\end{align}

The co-benefit sub-indicators take per-typology default levels, each cited in
the typology library (\Cref{app:typologies}), which the user may override.

\paragraph{The sub-indicator rules, fixed at version \methver{}.} An earlier
version of this document disclosed that the rules mapping user inputs to the
individual sub-indicator scores were unspecified and would be fixed alongside
the engine implementation. They are now fixed, live in
\texttt{config/derived\_scores.yaml}, and --- like the aggregation weights ---
are declared expert judgment, versioned and adjustable rather than derived from
literature. All suitability sub-indicators derive from inputs the user has
already supplied; no new questions are asked.

\begin{table}[H]
\centering
\caption[Suitability sub-indicator rules]{Derivation of the NbS Suitability
sub-indicators (decision D-022). A score of 25 on space, soil, or water
coincides with the hard ``not suitable'' flag of \Cref{sec:suitability};
an unknown availability never raises a flag, because a disqualification is
never asserted from absent information.}
\label{tab:suitability-rules}
\small
\begin{tabularx}{\textwidth}{@{}l>{\raggedright\arraybackslash}X@{}}
\toprule
\textbf{Sub-indicator} & \textbf{Rule} \\
\midrule
Space &
$r = \text{site area} / \text{typology minimum viable area}$:
$r < 1$ $\rightarrow$ 25 \emph{and the flag}; $1 \le r < 2$ $\rightarrow$ 50;
$2 \le r < 5$ $\rightarrow$ 75; $r \ge 5$ $\rightarrow$ 100. \\
\addlinespace
Soil &
Availability against the typology requirement on the ordinal scale
\emph{none} $<$ \emph{limited} $<$ \emph{moderate} $<$ \emph{high}:
requirement \emph{none} $\rightarrow$ 100; availability unknown $\rightarrow$
50; below requirement $\rightarrow$ 25 \emph{and the flag}; meets exactly
$\rightarrow$ 75; exceeds $\rightarrow$ 100. \\
\addlinespace
Water &
The same rule on \emph{none} $<$ \emph{occasional} $<$ \emph{reliable} against
the irrigation requirement. \\
\addlinespace
Maintenance &
Declared maintenance intensity through the inverted scale: low $\rightarrow$
100, medium $\rightarrow$ 50, high $\rightarrow$ 25, unknown $\rightarrow$ 50. \\
\addlinespace
Urban context &
Site land use inside the typology's typical-use contexts $\rightarrow$ 100;
outside them $\rightarrow$ 25 (a caution, not a disqualification); unknown
$\rightarrow$ 50. \\
\bottomrule
\end{tabularx}
\end{table}

The four equity sub-indicators read how much equity benefit the intervention
can deliver at this site, so a \emph{deficit} raises the score, paralleling the
risk framing of the Vulnerability Score: vulnerable-user benefit reuses the
vulnerable-population input already collected; public accessibility and
community participation are the two new optional inputs introduced at this
version; and safety concern deliberately uses the standard (non-inverted)
scale, because a site with greater safety and comfort problems stands to gain
more from a well-designed intervention.

Two disclosures carry over unchanged from the previous version. The co-benefit
weights include a \emph{social inclusion} term for which the typology library
supplies no default; absent a user override it falls to the neutral default of
\Cref{sec:unknown} and is itemised as an applied assumption. And the Equity
Score of \Cref{eq:equity} is computed and reported as an output block with its
own confidence rating but does \emph{not} enter the final aggregation of
\Cref{sec:aggregation}: equity influences the headline score through the
Vulnerability Score instead. Whether that is the right design remains a
legitimate question for reviewers.

\subsection{Energy savings --- derived, not assumed}
\label{sec:energy}

An earlier draft of the methodology stored per-typology ``energy reduction
factors'' in the range \SIrange{2}{15}{\percent}. No source could be found for
any of them. They were removed and replaced by a derivation from the tool's own
estimated cooling effect:

\begin{equation}
\label{eq:energy}
E \;=\; D \times \dT \times s,
\qquad s \in [0.02,\, 0.04]\ \text{per \si{\degreeCelsius}},
\end{equation}

where $D$ is the user-supplied annual cooling energy demand in
\unit{\kilo\watt\hour} and $s$ is the temperature--electricity-demand
sensitivity. Because both $\dT$ and $s$ are intervals, the result is an
interval:

\begin{equation}
\label{eq:energy-range}
E_{\min} = D \times \dT_{\min} \times 0.02,
\qquad
E_{\max} = D \times \dT_{\max} \times 0.04 .
\end{equation}

\paragraph{Basis.} \textcite{akbari2001} reports that urban electricity demand
increases by \SIrange{2}{4}{\percent} for each \qty{1}{\degreeCelsius} of
temperature increase, and estimates that \SIrange{5}{10}{\percent} of urban
electricity demand serves to compensate for the
\SIrange{0.5}{3.0}{\degreeCelsius} urban temperature elevation. Applying that
sensitivity in reverse to an estimated temperature reduction yields an energy
saving that is sourced, internally consistent with the tool's own cooling
estimate, and automatically responsive to site conditions --- a poorly suited
site yields a smaller $\dT$ and therefore a smaller energy saving, with no
separate parameter to maintain.

\paragraph{Preconditions.} Energy savings are calculated only when all three of
the following hold: the user confirms that nearby building cooling demand is
relevant to the intervention; an annual cooling energy demand is supplied; and
the selected typology is flagged as capable of affecting adjacent building
loads. Otherwise the result carries an explicit status, checked in a fixed order ---
\texttt{typology\_not\_applicable} (the intervention cannot affect building
loads, whatever the user supplied), \texttt{not\_applicable} (the user states
demand is not relevant), \texttt{relevance\_not\_confirmed} (the relevance
question was skipped or answered \emph{unknown}; a supplied demand figure
cannot substitute for the unanswered confirmation), or
\texttt{missing\_energy\_demand} --- and \textbf{never a zero}, which would
understate the result while appearing to be a finding.

\paragraph{An acknowledged weakness.} The sensitivity is drawn from
predominantly North American evidence, and its transferability to other
building stocks, construction practices, and climates is untested. Building
stock differs in envelope performance, in the prevalence and type of mechanical
cooling, and in occupant behaviour, all of which affect how a change in outdoor
air temperature translates into a change in electricity demand. This is an area
where reviewer input is specifically sought. It is also worth noting that
\textcite{akbari2001} is recorded in \Cref{app:verification} as
\emph{metadata + finding (secondary)}: the bibliographic record and the
\SIrange{2}{4}{\percent} figure were confirmed from multiple indexing records
and an institutional listing, not from a retrieved full text.

\subsection{Greenhouse gas emissions avoided}
\label{sec:ghg}

\begin{equation}
\label{eq:ghg}
G \;=\; E \times \phi,
\end{equation}

where $\phi$ is the grid emission factor in
\unit{\kilo\gram\per\kilo\watt\hour} of carbon dioxide equivalent, applied to
both ends of the energy interval.

Grid emission factors are sourced per country from \textcite{ember}, which
publishes electricity carbon intensity for 215 countries under a CC-BY-4.0
licence permitting redistribution with attribution inside an Apache-2.0 tool.
Where no verified factor is available for a country, the tool reports
greenhouse gas savings as \emph{not calculated} rather than substituting a
global average.

\begin{keynote}
\noindent\textbf{At version \methver{} no country emission factors ship with
the tool.} The country defaults file declares the source, licence, unit, and
per-entry schema, but its country table is empty pending per-country
verification. In consequence, a deployment using the shipped configuration
unmodified will report greenhouse gas savings as \emph{not calculated} for
every country. This is the intended failure mode of the verification policy ---
no unverified value ships --- but it means \Cref{eq:ghg} is at present a
specified formula awaiting its data, and any assessment reporting a greenhouse
gas figure is necessarily using a locally supplied factor.
\end{keynote}

\subsection{Costs --- why the tool ships no default values}
\label{sec:costs}

\textcite{worldbank2021} states that nature-based solution project costs vary
significantly and are highly site- and project-specific, that unit costs differ
sharply between developed and developing contexts --- illustrated by dredging
at \qty{2}{\USD\per\cubic\metre} in Bangladesh against
\qty{59}{\USD\per\cubic\metre} in the United Kingdom, roughly a thirtyfold
difference for the same operation --- and that urban implementation costs
typically exceed rural ones. No source was found offering globally applicable
unit costs for the typologies in this library.

The tool therefore ships \textbf{no default cost values} in version 1. Where a
user supplies a capital cost, the tool computes cost outputs; where they do
not, capital cost, payback, and cost feasibility are all reported as
\emph{not estimated}.

The reasoning is worth stating explicitly, because shipping a plausible default
would have been easy and would have made the tool look more complete.
Fabricating a default per-square-metre cost would generate the most
decision-relevant number in the entire assessment --- simple payback --- from
an invented input. No confidence rating adequately qualifies a payback figure
whose numerator was guessed: the reader sees a number of years, and the number
of years is what enters the decision. Locally calibrated cost tables remain a
defined extension point; a deployment may supply its own cost configuration,
which the tool records as a documented assumption in the result.

Where costs are supplied:

\begin{align}
\label{eq:annual-savings}
C_{\mathrm{annual}} &= E \times p, \\
\label{eq:payback}
T_{\mathrm{payback}} &= \frac{C_{\mathrm{capital}}}{C_{\mathrm{annual}}},
\qquad \text{undefined if } C_{\mathrm{annual}} \le 0,
\end{align}

with $p$ the energy price per \unit{\kilo\watt\hour}. Because $E$ is an
interval, so is the payback, with the bounds inverted: the largest energy
saving gives the shortest payback.

Cost feasibility is \emph{derived} rather than user-asserted:

\begin{equation}
\label{eq:cost-feasibility}
S_{\mathrm{cost}} \;=\; f\bigl(\text{payback bracket},\;
\text{implementation complexity},\; \text{maintenance intensity}\bigr),
\end{equation}

with short payback, low complexity, and low maintenance yielding high
feasibility. The bracket boundaries and combination rule, fixed at version
\methver{} (decision D-023), are as follows. Because the energy saving is an
interval spanning an order of magnitude, the payback bracket is applied to the
payback computed from the \emph{central} energy estimate --- the midpoint of
the savings range --- while the reported payback interval still carries both
ends. The brackets are half-open and follow common public-investment screening
horizons: under 5 years scores 100 (\emph{short}); 5--10 years 75
(\emph{medium}); 10--20 years 50 (\emph{long}); 20 years and beyond 25
(\emph{very long}). The combination is

\begin{equation}
\label{eq:cost-feasibility-rule}
S_{\mathrm{cost}} = 0.50\,S_{\mathrm{bracket}}
 + 0.25\,S_{\mathrm{complexity}} + 0.25\,S_{\mathrm{maintenance}},
\end{equation}

with complexity and maintenance normalised through the inverted scale so that
low complexity and low maintenance raise feasibility, and unknowns taking the
neutral 50, itemised. Payback carries half the weight as the only quantitative
economic signal. An \emph{investment readiness} level is reported alongside:
the bracket sets the base (\emph{short} $\rightarrow$ high, \emph{medium}
$\rightarrow$ medium, otherwise low), downgraded one level for high
implementation complexity and one level for low energy-block confidence, with a
floor at low; it is \emph{not estimated} whenever cost feasibility is.

\paragraph{Exclusion rather than substitution.} When payback cannot be
computed, cost feasibility is reported as unavailable and is \emph{excluded}
from the final aggregation, with the remaining weights redistributed
proportionally (\Cref{eq:aggregation-redistributed}). The alternative ---
substituting a neutral 50 --- would silently reward projects that have provided
no economic evidence at all, placing them ahead of projects that supplied a
genuinely poor payback figure. Exclusion preserves the ordering among options
that did supply data, and the missing block is reported rather than hidden.

\subsection{Aggregation and the weighting question}
\label{sec:aggregation}

The headline score is a weighted sum of six components:

\begin{equation}
\label{eq:aggregation}
S_{\mathrm{opp}} =
  0.25\,\mathrm{HPI}
+ 0.25\,S_{\mathrm{cool}}
+ 0.15\,S_{\mathrm{suit}}
+ 0.15\,S_{\mathrm{vuln}}
+ 0.10\,S_{\mathrm{cobenefit}}
+ 0.10\,S_{\mathrm{cost}} .
\end{equation}

Where a component is unavailable --- in practice, cost feasibility --- the sum
is taken over the available set $\mathcal{A}$ and renormalised:

\begin{equation}
\label{eq:aggregation-redistributed}
S_{\mathrm{opp}} \;=\;
\frac{\displaystyle\sum_{i \in \mathcal{A}} w_i\,S_i}
     {\displaystyle\sum_{i \in \mathcal{A}} w_i} .
\end{equation}

The result is interpreted through the bands of \Cref{tab:hpi-bands}.

\paragraph{These weights are expert judgment, and the tool says so.}
\textcite{nardo2008} is explicit that weighting in composite indicators is a
value choice rather than a technical derivation, and no literature can supply
the relative importance of cooling performance against equity against cost. The
weights are declared here, versioned in configuration, adjustable by any
deployment, and are to be defended empirically through the sensitivity analysis
specified in \Cref{sec:sensitivity} rather than by appeal to authority.

\subsection{The equity-forward weighting, stated openly}
\label{sec:effective-weights}

Vulnerability enters the final score \emph{twice}. It contributes \num{0.40} of
the Heat Priority Index, which itself carries \num{0.25} in
\Cref{eq:aggregation}, and it appears again directly at \num{0.15}. Decomposing
the Heat Priority Index in \Cref{eq:aggregation} gives

\begin{align}
S_{\mathrm{opp}} &=
  0.25\bigl(0.60\,S_{\mathrm{heat}} + 0.40\,S_{\mathrm{vuln}}\bigr)
+ 0.25\,S_{\mathrm{cool}} + 0.15\,S_{\mathrm{suit}} + 0.15\,S_{\mathrm{vuln}}
+ 0.10\,S_{\mathrm{cobenefit}} + 0.10\,S_{\mathrm{cost}} \notag \\[2pt]
\label{eq:effective}
&= 0.15\,S_{\mathrm{heat}} + 0.25\,S_{\mathrm{vuln}} + 0.25\,S_{\mathrm{cool}}
 + 0.15\,S_{\mathrm{suit}} + 0.10\,S_{\mathrm{cobenefit}}
 + 0.10\,S_{\mathrm{cost}} .
\end{align}

The effective weights are set out in \Cref{tab:effective-weights}.

\begin{table}[H]
\centering
\caption[Effective weights in the final score]{Effective weights in the NbS
Cooling Opportunity Score, after decomposing the Heat Priority Index. The
nominal weights are those of \Cref{eq:aggregation}; the effective weights are
those of \Cref{eq:effective}. Both columns sum to unity, as the total row
records. The two rows below that total are not additional weights: they
restate, for emphasis, the totals that vulnerability and heat exposure reach
once the Heat Priority Index is decomposed.}
\label{tab:effective-weights}
\small
\begin{tabular}{@{}lrr@{}}
\toprule
\textbf{Component} & \textbf{Nominal weight} & \textbf{Effective weight} \\
\midrule
Heat Priority Index & 0.25 & --- \\
\quad of which heat exposure   & & 0.15 \\
\quad of which vulnerability   & & 0.10 \\
Cooling Potential   & 0.25 & 0.25 \\
NbS Suitability     & 0.15 & 0.15 \\
Vulnerability (direct) & 0.15 & 0.15 \\
Co-benefits         & 0.10 & 0.10 \\
Cost Feasibility    & 0.10 & 0.10 \\
\midrule
\textbf{Total}      & \textbf{1.00} & \textbf{1.00} \\
\midrule
\multicolumn{3}{@{}l}{\itshape Restated from the rows above, not added to them:}\\
\textbf{Vulnerability, total} & & \textbf{0.25} \\
\textbf{Heat exposure, total} & & \textbf{0.15} \\
\bottomrule
\end{tabular}
\end{table}

\begin{keynote}
\noindent The tool therefore weights \emph{who is affected} (\num{0.25}) above
\emph{how hot the site is physically} (\num{0.15}). Between two equally hot
sites, the one serving a more vulnerable population ranks higher.

\textbf{This is a value position, not a technical necessity.} There is no
finding in the literature that requires it, and a reasonable reviewer may
reject it. A city whose mandate is to reduce peak thermal load across its
territory, rather than to direct adaptation investment toward disadvantaged
communities, would be right to rebalance it.

It is disclosed here, and quantified, for a specific reason: an undisclosed
double contribution would be a defect, indistinguishable from an arithmetic
error, and a reviewer discovering it unaided would be entitled to distrust
everything else in the methodology. Disclosed and quantified, it is a
defensible position that reviewers may legitimately dispute. Any deployment may
rebalance it in \texttt{config/weights.yaml} without touching code, and every
result is to record the methodology version under which it was produced.
\end{keynote}

The double contribution is also not an accident of construction that happened
to be noticed afterwards. It follows from a deliberate decision, recorded in
the project's decision log as D-007, that the tool should prioritise who is
affected slightly above how hot a place is. The alternative designs --- dropping
the direct vulnerability term, or removing vulnerability from the Heat Priority
Index --- were both available; each would have produced a defensible instrument
with a different value stance.
